\section{Introduction}

Canadian Fish --- Literature in most of the published literature --- is a
six-player, two-team card game of imperfect information. On your turn you name a
card and an opponent; if they hold it you take it and continue, and if they do
not, the turn is theirs. At any moment either team may claim a half-suit by
naming which teammate holds each of its six cards, and any error in that claim
hands the half-suit to the opponents. Everything that happens is public.

The game has almost no computational literature. There is no arXiv paper on it,
no engine, and one open-source learning agent that plays a different variant.
That absence is not because the game is easy: it has six players, two teams that
cannot communicate, an information set of astronomical size, and a horizon of
sixty to a hundred and twenty decisions. It is closer to Bridge or Skat than to
anything solved.

FishBot v0.3 \cite{fishbot03} was the first serious attempt: a transparent
one-ply utility policy over an approximate posterior, fitted on 126{,}600
simulated games and evaluated on held-out seeds. It established that remembering
public asks and reconciling them against card counts dominate everything else,
and it left three things open --- an approximate belief model, a confidence
threshold in place of a declaration policy, and a simulator that did not
implement the rule allowing a claim at any moment.

This paper closes all three. The route runs through one structural observation.

\begin{quote}
\emph{Cards in Fish move only through publicly observed transfers. Therefore a
card whose location has never been revealed is still with whoever was dealt it,
and the entire hidden state of the game is the initial deal.}
\end{quote}

The consequences are unusually strong. Inference becomes a degree-constrained
counting problem over a capacity vector, which we solve exactly --- exact
marginals, exact probabilities for a named allocation, and rejection-free
sampling --- for a few hundred thousand arithmetic operations
(\S\ref{sec:belief}). The one constraint that resists a product form, the
certificate that an ask carries about the asker's own hand, is confined to a
single half-suit and can therefore be enumerated inside it exactly rather than
approximated. And a one-line reading of ask legality yields a strategic result
that no previous agent implements: a half-suit held entirely by one team can
never be asked in by the other, so it can never be taken back, so waiting to
claim it is free (\S\ref{sec:theory}). Declaration stops being a threshold and
becomes an optimal-stopping problem.

\subsection*{Claims}

This paper defends four bounded claims.

\begin{enumerate}
  \item \textbf{The exact posterior is computable and is computed.} Belief
        marginals, the probability of a named allocation, and consistent deal
        samples are exact with respect to every deduction the rules license, and
        the implementation is validated against brute force, against an
        independent dynamic program, and against exact rejection sampling.
  \item \textbf{A half-suit held entirely by one team cannot be taken back}
        (Theorem~\ref{thm:locked}), so declaration is optimal stopping rather
        than a threshold. The agent that follows declares near-perfectly; we are
        explicit in \S\ref{sec:ablations} that the ablations credit the exact
        allocation probability rather than the stopping rule itself.
  \item \textbf{FishBot v0.4 is the strongest policy in the tested population},
        and --- the point of the exercise --- it is strongest against
        \emph{every} member of that population, not merely on average.
  \item \textbf{The forecasts it stops on are calibrated}, which is what makes
        the third claim a statement about decision quality rather than about
        one seed bank.
  \item \textbf{It is much harder to counter than its predecessor.} An adversary
        fitted specifically to beat it, in the same policy class and with the
        same optimiser and budget, does not reach 50\% against it; the identical
        procedure beats FishBot v0.3 comfortably. This is a lower bound on
        exploitability rather than a bound, and we say so, but the comparison
        between the two figures is the meaningful part.
\end{enumerate}

A sixth result was not among the claims we set out to make and is reported
because it is a fact about the game rather than about the agent: because a
half-suit frozen by Theorem~\ref{thm:locked} also admits no further information
about itself, two policies that both correctly decline to cash an uncertain one
deadlock permanently, and the rules as usually stated therefore do not guarantee
termination against deterministic play (\S\ref{sec:termination}).

This paper does not claim a Nash or team equilibrium, an exploitability bound, or
that Fish is solved. \S\ref{sec:limitations} is explicit about what remains, including
the one caveat that most constrains the word ``exact'': the posterior is exact
with respect to the \emph{rules}, and a fully Bayesian agent would additionally
weight deals by how likely the observed actions are under each opponent's
particular playstyle. We implement that correction, fit it, and report that it
adds nothing once the rule-based certificates are enforced --- which is evidence
about this population rather than a general theorem.
